\documentclass[11pt,a4paper]{article}
\usepackage[margin=2.5cm]{geometry}
\usepackage[T1]{fontenc}
\usepackage{booktabs}
\usepackage{xcolor}
\usepackage{listings}
\usepackage[hidelinks]{hyperref}

\emergencystretch=2em
\definecolor{codebg}{gray}{0.95}
\lstset{language=C++, basicstyle=\ttfamily\small, keywordstyle=\color{blue!70!black},
  commentstyle=\color{green!50!black}\itshape, backgroundcolor=\color{codebg},
  frame=single, rulecolor=\color{black!20}, breaklines=true, showstringspaces=false,
  tabsize=2, columns=flexible}

\title{ATmega32A Module \\ \Large Common Definitions (\texttt{common.hpp})}
\author{Ali Sahafi}
\date{}

\begin{document}
\maketitle

\section{What this file is}

Every driver module in this course (GPIO, ADC, UART, \dots) starts with
\lstinline|#include "../common/common.hpp"|. You never include this file
yourself --- it comes along automatically with whichever module you use.
It provides three things:

\begin{enumerate}
  \item The MCU definition \texttt{\_\_AVR\_ATmega32A\_\_} so that
        \texttt{<avr/io.h>} exposes the correct register names.
  \item The standard AVR headers: \texttt{<avr/io.h>} (registers),
        \texttt{<avr/interrupt.h>} (ISR support, \texttt{sei()}/\texttt{cli()}) and
        \texttt{<util/delay.h>} (\texttt{\_delay\_ms()}).
  \item A set of Arduino-like constants used across all modules.
\end{enumerate}

\section{Constants}

\begin{center}
\begin{tabular}{lll}
\toprule
\textbf{Constant} & \textbf{Value} & \textbf{Meaning} \\
\midrule
\texttt{INPUT}        & 0    & Pin is an input (high impedance) \\
\texttt{OUTPUT}       & 1    & Pin is an output \\
\texttt{INPUT\_PULLUP} & 2    & Input with the internal pull-up resistor enabled \\
\texttt{LOW}          & 0    & Logic 0 (GND) \\
\texttt{HIGH}         & 1    & Logic 1 (VCC) \\
\texttt{ALL}          & 0xFF & Selects the \emph{whole port} instead of one pin \\
\bottomrule
\end{tabular}
\end{center}

\section{The \texttt{F\_CPU} rule}

\texttt{F\_CPU} tells the toolchain how fast your CPU clock is. Delay loops,
UART baud rates and timer calculations are all derived from it, so it must
match the \emph{real} clock of your chip (set by the fuse bits --- see the
\texttt{Makefile} targets \texttt{int-1mhz} \dots\ \texttt{ext-xtal}).

\textbf{Always define it \emph{before} the first driver include:}

\begin{lstlisting}
#define F_CPU 8000000UL   // 8 MHz -- must come first!
#include "drivers/gpio/gpio.hpp"
\end{lstlisting}

If you forget, the driver falls back to 8\,MHz and the compiler prints a
warning. A wrong \texttt{F\_CPU} does not stop compilation --- but your
delays will be wrong and your UART will print garbage. If the serial
terminal shows random characters, check \texttt{F\_CPU} and the fuses first.

\section{One source file per program}

The drivers are \emph{header-only}: all code lives in \texttt{.hpp} files and
each firmware image is built from exactly one \texttt{.cpp} file. Include
only the modules you need --- unused peripherals then cost no flash memory
and no interrupt vectors.

\end{document}
